\chapter{SYSTEM ANALYSIS}
\section{Scope of the Project}
This project aims to develop a web-based platform that connects passengers and administrators through a centralized bus booking system. The system allows:
\begin{itemize}
    \item Passengers to register, log in, search for route availability, view schedules, select seats interactively, and process simulated payments.
    \item Administrators to manage bus fleets, update routes, adjust fares, and monitor total capacity.
\end{itemize}

\section{Overall Description}
\textbf{Product Perspective}

The Bus Ticket Booking System is a web-based application designed as a standalone system that operates on a client-server architecture. The system is developed using the MERN Stack, where the frontend is built using React.js, the backend logic is handled by Node.js and Express.js, and data storage is managed using MongoDB.

The application acts as a centralized platform that bridges the gap between passengers and transport operators by providing a digital solution for booking, seat management, and ticket tracking.
The system does not replace existing physical buses entirely but enhances them by automating ticket allocations, double-booking prevention, and scheduling.

\textbf{Product Functions}

The major functions of the system are categorized based on user roles:

1. Passenger Functions
\begin{itemize}
    \item User registration and secure login
    \item Bus search using filters such as date, source, and destination
    \item Interactive seat selection mapped to bus layouts
    \item Passenger detail entry and simulated checkout
    \item Viewing booking history and cancellation
\end{itemize}

2. Administrator Functions
\begin{itemize}
    \item Managing bus schedules and fleet information
    \item Monitoring active bookings to prevent race conditions
    \item Adjusting fare pricing structures
\end{itemize}

\section{Project Pipeline}

\begin{figure}[H]
\centering
\includegraphics[width=0.8\textwidth]{project_pipe.png} 
\caption{Project Pipeline}
\end{figure}

\subsection{Software Requirement Specification (SRS)}

\begin{table}[H]
\centering
\begin{tabular}{|p{1cm}|p{3cm}|p{2cm}|p{6cm}|p{3.5cm}|}
\hline
\textbf{No.} & \textbf{Requirement} & \textbf{Essential / Desirable} & \textbf{Description} & \textbf{Remarks} \\ \hline
RS1 & User Registration \& Login & Essential & Users register using email and login securely to access checkout. & Restricts booking to logged-in users. \\ \hline
RS2 & Bus Search Engine & Essential & Users can search buses by source, destination, and exact dates. & Filters available routes dynamically. \\ \hline
RS3 & Seat Selection Map & Essential & Visual interface for users to pick available seats only. & Real-time locking prevents double booking. \\ \hline
RS4 & Passenger Details & Essential & Collecting mandatory information per selected seat. & Ensures ticket validity. \\ \hline
RS5 & Booking History & Essential & Profile section displaying past and upcoming journeys. & Allows for cancellations. \\ \hline
RS6 & Admin Dashboard & Essential & View and manage all bus schedules and base fares. & Admin has full control over system. \\ \hline
\end{tabular}
\caption{Software Requirement Specification}
\end{table}

\textbf{Assumptions}

\begin{table}[H]
\centering
\begin{tabular}{|p{1cm}|p{14cm}|}
\hline
\textbf{No.} & \textbf{Assumption} \\ \hline
AS1 & Users have a valid email address mapped to their account. \\ \hline
AS2 & Internet connectivity is available for all users accessing the booking engine. \\ \hline
AS3 & Route details posted by administrators are accurate and authentic. \\ \hline
AS4 & Only verified administrators can alter the global seat pricing model. \\ \hline
AS5 & A payment gateway simulation reflects genuine checkout scenarios. \\ \hline
\end{tabular}
\caption{System Assumptions}
\end{table}

\section{User Classes and Characteristics}
The system identifies two main user classes:
\begin{itemize}
    \item \textbf{Passengers:} Individuals searching for travel tickets. They possess basic web application usage knowledge and interact with the system to book and cancel trips.
    \item \textbf{Administrator:} The admin has advanced system access and technical understanding, responsible for maintaining route integrity.
\end{itemize}

\section{Feasibility Study}
\textbf{1. Technical Feasibility}
The project is technically feasible as it uses widely adopted frontend technologies such as HTML, CSS, JavaScript, and React.js. The backend utilizes Express.js and MongoDB.
\textbf{2. Economic Feasibility}
The development of the frontend application does not involve significant cost, as open-source libraries are used.
\textbf{3. Operational Feasibility}
The system is easy to operate. Since the application is designed with intuitive navigation, users can book tickets without extensive training.

\section{Specific Requirements}
1. User Registration and Authentication
\begin{itemize}
    \item Passengers register securely using valid credentials.
    \item Upon login, sessions are maintained securely to verify checkout access.
\end{itemize}
2. Interactive Booking System
\begin{itemize}
    \item The system visually displays available, occupied, and selected seats on a live matrix layout.
\end{itemize}
3. Search and Filtering
\begin{itemize}
    \item Direct searching between source and destination endpoints based on calendar dates.
\end{itemize}
4. Booking Integrity
\begin{itemize}
    \item Application validates concurrent seat selections via backend locks ensuring no double-booking accidents.
\end{itemize}

\section{External Interface Requirements}
\textbf{1. User Interface}
The Bus Ticket Booking Portal provides a responsive web interface accessible through standard browsers, presenting dashboards, search forms, and seat grids. 
\textbf{2. Hardware Interface}
Operates on desktops, laptops, tablets, and smartphones supporting a modern web browser.
\textbf{3. Software / Communication Interface}
Communicates over standard HTTP/HTTPS protocols with a MongoDB storage backend utilizing RESTful APIs for fetching dynamic prices and seat statuses.

\section{System Environment}
\subsection{Frontend Technology}
The MEARN stack is adopted for the project, with emphasis on the frontend components using HTML for structure, CSS for styling, JavaScript for logic, and React.js for interactive user interfaces.
\subsection{Backend Technology}
Node.js and Express.js handle API routing and concurrency requests, while MongoDB processes seat occupation records. JSON Web Tokens (JWT) act as the authentication middleware covering security restrictions across passenger operations.
